\documentclass[11pt,a4paper]{article}
\usepackage[T1]{fontenc}
\usepackage[utf8]{inputenc}
\usepackage{lmodern,microtype}
\usepackage[margin=.85in]{geometry}
\usepackage{amsmath,amssymb,amsthm,mathtools,aliascnt}
\usepackage{booktabs,tabularx,array,enumitem,needspace,float}
\usepackage[section]{placeins}
\usepackage{xcolor,tikz,mdframed}
\usetikzlibrary{arrows.meta,positioning,calc}
\usepackage[hidelinks]{hyperref}
\usepackage[nameinlink,noabbrev]{cleveref}
\usepackage{fancyhdr}
\hypersetup{pdftitle={Sync HotStuff RR v4.1},pdfsubject={Regional finality and L1-authorized external truncation}}
\setlength{\headheight}{14pt}
\pagestyle{fancy}\fancyhf{}
\lhead{\small Sync HotStuff RR v4.1}\rhead{\small L1-authorized regional sessions}
\cfoot{\thepage}
\setlength{\parindent}{0pt}\setlength{\parskip}{4pt}
\setlength{\emergencystretch}{2em}
\setlist{nosep,leftmargin=*}
\newcolumntype{Y}{>{\raggedright\arraybackslash}X}
\definecolor{rrRed}{HTML}{A51C30}
\definecolor{rrBlue}{HTML}{2563EB}
\definecolor{rrTeal}{HTML}{0F766E}
\newcommand{\RR}[1]{{\color{rrRed}#1}}
\newcommand{\QC}{\mathcal C}
\newcommand{\FC}{\mathcal F}
\newcommand{\Hash}{\mathsf H}
\newcommand{\rank}{\operatorname{rank}}
\newcommand{\wt}{\operatorname{wt}}
\newcommand{\tagg}[1]{\ensuremath{\mathsf{#1}}}
\newtheorem{definition}{Definition}[section]
\newaliascnt{assumption}{definition}
\newtheorem{assumption}[assumption]{Assumption}
\aliascntresetthe{assumption}
\crefname{assumption}{assumption}{assumptions}
\Crefname{assumption}{Assumption}{Assumptions}
\newcounter{protocol}
\newenvironment{protocol}[2]{\refstepcounter{protocol}\label{#1}
\begin{mdframed}[nobreak=true,linewidth=.6pt,skipabove=8pt,skipbelow=8pt,
innertopmargin=7pt,innerbottommargin=7pt]\small
\textbf{Protocol \theprotocol. #2}\par\smallskip}{\end{mdframed}}
\crefname{protocol}{Protocol}{Protocols}
\Crefname{protocol}{Protocol}{Protocols}
\title{Sync HotStuff RR v4.1\\\large Regional Finality and L1-Authorized Rotation}
\author{}\date{12 October 2026}
\begin{document}
\maketitle
\begin{abstract}
L1 supplies stake configurations, leader schedules, regional quotas, and
status updates. Within each region, Sync HotStuff commits DATA blocks
using proposal timers and rotates leaders on terminal certificates.
A single DRAIN position is available only when recovery opens on a
full-quota chain. Cross-region replicas continuously fetch regional
blocks and finalized L1 control state. Normal completion retains the
terminal DATA prefix. An activated force instruction instead fixes an
external cutoff chosen by an honest multisignature service. Old regional
blocks above that cutoff may remain locally committed but are excluded
from the external canonical history. Both paths use one bootstrap
handler. This draft specifies the model and protocols; security and
complexity analyses are deferred.
\end{abstract}

\section{Model and L1 control state}
\label{sec:model}
L1 is the public source of sequencing authority. It determines who may
vote and lead, how many DATA blocks a region may add, which regions are
eligible, and when an explicit forced rotation becomes effective.
Regional Sync HotStuff~\cite{sync} orders and commits blocks locally.
L1 force instructions determine which of those blocks remain in the
external canonical history.
Cross-region infrastructure distributes their contents, certificates, and
L1 proofs; it does not elect a leader by reporting a tip.

\subsection{What nodes read from L1}
\label{sec:l1}
A regional session $r$ is one turn of a physical region $i_r$.
Within that turn, $g$ counts recoveries and $t$ counts leader terms.
Nodes read five things from finalized L1:

\begin{table}[htbp]\centering\small
\caption{L1 information used by the protocol.}
\label{tab:l1}
\begin{tabularx}{\linewidth}{@{}p{.25\linewidth}Y@{}}
\toprule
Information & Meaning\\
\midrule
Stake configuration & The members, public keys, and voting weights of
each region, together with the protocol parameters.\\
Leader schedule & $L(r,g,t)=(i_r,\ell_{r,g,t})$: the region and leader
assigned to that view.\\
Regional quota & $TB_r$: the number of new DATA blocks allowed in
session $r$. Local recovery does not reset it.\\
Regional status & $\chi(i)\in\{\tagg{live},\tagg{off}\}$: whether a region
is eligible for a future turn.\\
Force instruction & $\tagg{FORCE}(r,h_*)$, with an activation height
$a$: end session $r$'s external history at DATA height $h_*$ once L1
through $a$ is finalized. The instruction contains no anchor block hash.\\
\bottomrule
\end{tabularx}
\end{table}

The multisignature service posts status updates to L1. Write update $j$ as
\[
 E_j=(a_j,\chi_j,\phi_j),\qquad
 \phi_j\in\{\bot,\tagg{FORCE}(r_j,h_j^*)\}.
\]
The symbol $\bot$ means that the update changes status only.
A force instruction also marks its target region off for future selection.
Announcement permits preparation; activation changes the external
authorization. Neither event is a regional BFT quit proof.

Nodes verify the update's authorization and L1 finality. Updates have
consecutive versions and increasing activation heights. Once finalized,
an update cannot be rewritten. If it was posted at L1 height $u_j$, its
activation satisfies $a_j\ge u_j+w_{\min}$ for the configured notice
distance $w_{\min}>0$. At finalized control height $\lambda$,
\[
 j(\lambda)=\max\{j:a_j\le\lambda\},\qquad
 \chi_{\rm eff}(\lambda)=\chi_{j(\lambda)}.
\]
The initial entry $E_0=(a_0,\chi_0,\bot)$ supplies the initial map.
All control heights are at least $a_0$.

We write $\tagg{AUTH}_{r,g}$ for the configuration, leader schedule,
and quota used by generation $(r,g)$. These may be derived from preset
L1 rules; a new L1 transaction is not needed for every local rotation.
The initial region is the first live region in the configured cyclic
list $\mathcal L_R$. It opens at genesis with a synthetic regional
receipt. If no region is eligible, initialization waits.

\begin{assumption}[Timely L1 authority and common control order]
\label{ass:authority}
For every view the protocol may enter, including a pending recovery view,
all correct members have its matching finalized authority before its
first permitted entry. L1 progresses independently of regional BFT.
Control updates and completed handoffs have one authenticated common
order. This order fixes one successor and one retained DATA prefix at
each sealed boundary. Later updates cannot change a sealed boundary.
\end{assumption}
A boundary is \emph{sealed} when this common order fixes its anchor,
successor, and selection height. We retain this ordered-control service
as an external assumption, including when block production is silent.
Local message arrival order does not decide between competing handoffs.
The protocols do not add a separate L1 transaction for each regular
handoff.

\subsection{Stake, communication, and recovery services}
For region $i$, let $W_i=\sum_{v\in V_i}\omega_{i,v}$ with
$\omega_{i,v}>0$. Byzantine stake is at most $F_i<W_i/2$.
Every regional certificate in this draft requires distinct signers of
total weight strictly greater than $F_i$. Equivalently, for a fault-share
bound $f_i=F_i/W_i<1/2$, its weight exceeds $f_iW_i$.
Membership and weights are fixed throughout a session, including local
recoveries. Slashing or a passive status entry does not change its
denominator. A later session verifies old evidence against the old snapshot.
The bound applies to the cumulative compromised signing weight of each
snapshot for as long as its certificates are accepted; rotating corruption
or later compromise of more historical keys is outside this model.

\begin{assumption}[Authenticated regional synchrony and global clock]
\label{ass:network}
While a regional session operates under the normal model, every
correct-to-correct message, including its validation data and processing,
arrives within $\Delta_i$; actual delay may be smaller. Replicas share a
global clock, retain their protocol state, and have unforgeable signatures
and binding hashes. Correct members of a prospective successor also
receive and validate its bootstrap bundles while awaiting activation.
\end{assumption}
The fault bound and synchrony condition are per region, not merely bounds
on the aggregate network. Cross-region catch-up and L1 dissemination have
their own service delays. They are not included silently in $\Delta_i$.

\begin{assumption}[Cross-region retrieval and bootstrap readiness]
\label{ass:infra}
Prospective successors continuously receive or fetch leading proposals,
QCs, regional finality receipts, ancestry, and finalized L1 control proofs.
The infrastructure provides leader push and bounded-range poll--fetch
repair from retained holders. For every pending authorized handoff with
a successor satisfying the regional service assumptions, the complete
entry bundle and required validation state eventually become available.
The readiness service delivers them to that successor's correct members
with entry skew at most $\Delta_i$. Each first recipient relays the bundle
within its region.
\end{assumption}
Cross-region catch-up may take longer than $\Delta_i$ before admission.
An unconfirmed tip is a retrieval hint; a fetch timeout grants no
production authority. Correct nodes retain the data and request history
needed to reconstruct the selected anchor.

\begin{assumption}[Honest multisignature service and external truncation]
\label{ass:cut}
The multisignature service is honest. A force instruction selects a
unique DATA block $B_{h_*}$ in a prefix through which the regional BFT
operated correctly and committed its blocks. That block, its regional
receipt, and its ancestry are recoverable outside the old region.
On activation, finalized L1 fixes session $r$'s external endpoint at
$h_*$, without changing an earlier sealed prefix. The trusted control
service continues to give functioning regions enough uninterrupted
service to make retained DATA progress and complete handoffs; adversarial force churn
is outside this model.
\end{assumption}
L1 is authoritative for this external cutoff. The instruction need not
cover every regional commitment: the old region may already have committed,
or may later commit, blocks above $h_*$. Those blocks are excluded from
the external canonical history. No proof of their absence, old-region
shutdown, or new acknowledgment from the failed region is required.
The regional receipt determines the anchor hash; the multisignature
instruction specifies only its height and session.

A force does not cancel regional timers, change local locks, reduce
$TB_r$, or undo the old region's committed log. It acts on the external
history assembled from regional outputs. The scope of irrevocable
external finality is defined in \cref{sec:external}.

\paragraph{Leader service and pending requests.}
The schedule gives correct leaders repeated opportunities in
\emph{actually entered} views, including repeated recovery initializers.
If DATA production continues across sessions, correct leaders receive
repeated actual DATA-producing turns; assignments to unreachable future
terms do not satisfy this condition. Exposed equivocating leaders are
eventually excluded by future authorities. These are service assumptions,
not consequences of stake deposits.

Correct-client requests eventually reach the serving regions and the
cross-region pending pool. They have a common admission order, finitely
many predecessors, and each request fits in a batch. Correct leaders
include the oldest admissible pending requests that fit and are not
already in their proposed ancestry. Empty DATA is allowed when no such
request is available. The external pool removes a request only when it
belongs to a sealed retained prefix. Requests in a force-discarded suffix
remain pending and are rechecked for inclusion in the successor.
Regional commitment alone does not remove them from this external pool.

\section{Views, blocks, and admission}
\label{sec:definitions}
A global view is $(r,g,t)$, ordered lexicographically. Here $r$ counts
regional sessions, $g$ counts recoveries within one session, and $t$ counts
ordinary leader terms within one generation. Each newly opened region
starts at $(r,0,0)$.

\begin{table}[htbp]\centering\small
\caption{View changes and their admission events.}
\label{tab:views}
\begin{tabularx}{\linewidth}{@{}p{.24\linewidth}p{.27\linewidth}Y@{}}
\toprule
Event & Target & Admission and timer treatment\\
\midrule
In-region normal rotation & $(r,g,t+1)$ & A term-terminal DATA QC below
the regional quota. Existing proposal timers continue.\\
In-region Quit & $(r,g+1,0)$ & A complete quit proof, the recovery wait,
and the next authority. Unfinished old-generation timers are canceled.\\
Regional quota completion & $(r+1,0,0)$ & A finalized terminal DATA anchor
and a sealed regular-handoff witness naming the successor.\\
Activated force & $(r+1,0,0)$ & Finalized L1 activation fixes an external cutoff; a regional
receipt identifies the anchor. Old regional timers are unaffected.\\
\bottomrule
\end{tabularx}
\end{table}
A DRAIN QC causes no view change. A status update without an applicable
force tag does not increment $r$, $g$, or $t$. Local recovery never resets
the regional quota.

\subsection{Two positions and the regional quota}
\label{sec:positions}
The unchanged DATA block $A_r$ is the session's cross-region anchor;
$A_0$ is genesis. Let $h(B)$ be global DATA height and $k_r(B)$ the
physical position in session $r$. Assign $k_r(A_r)=0$ for this session's
arithmetic, without rewriting the anchor's original signed fields.
A fresh physical block has the form
\[
 B=(r,k,h,\mathit{type},x,\Hash(B_{\rm parent}),\rho),
 \qquad \mathit{type}\in\{\tagg{DATA},\tagg{DRAIN}\},
\]
where $k=k_r(B_{\rm parent})+1$. DATA contains the next application batch
and has $h=h(B_{\rm parent})+1$. DRAIN has empty application contents
and $h=h(B_{\rm parent})$. Define
\[
 H_r=h(A_r)+TB_r.
\]
DATA ancestry may contain at most $TB_r$ fresh DATA blocks, including
uncommitted ones. An ANCHOR re-certification, a retransmission, or a DRAIN
does not consume this quota. The terminal DATA block $T_r$ has height $H_r$.

Each DATA block binds a finalized L1 reference $\rho$ with height
$\lambda(B)$, canonical hash, and registry proof. These heights do not
decrease along DATA ancestry; the first fresh DATA reference is at least
the session's opening control height. Regional proposal validation checks
the session's fixed configuration and leader authority. Later status or
force updates do not rewrite that local validity test.
DRAIN preserves its terminal DATA ancestor's reference.
External acceptance separately checks the applicable ordered control
history. Carrying an old L1 reference cannot bypass an activated cutoff.

\paragraph{Ordinary and shortened terms.}
Generation $(r,g)$ has an immutable opening $A_{r,g}$; $A_{r,0}=A_r$.
If $h(A_{r,g})<H_r$, term $t$ has
\[
 m_{r,g,t}=\min\{K,H_r-h(A_{r,g})-tK\}>0
\]
fresh DATA positions $j=0,\ldots,m_{r,g,t}-1$. Their positions are
\[
 k=k_r(A_{r,g})+tK+j+1,\qquad
 h=h(A_{r,g})+tK+j+1.
\]
There is no next term once its terminal DATA reaches $H_r$.
For example, $K=4$ and $TB_r=10$ give term lengths $4,4,2$ in an
uninterrupted generation starting at $A_r$.

\paragraph{The one permitted DRAIN position.}
\label{sec:drain}
If $g>0$ and $h(A_{r,g})=H_r$, this is a \emph{drain generation}.
Its only fresh position is $(t,j)=(0,0)$: one DRAIN extending $A_{r,g}$
with its current ANCHOR QC. The opening may itself be an older DRAIN.
There is no second fresh position or ordinary next term in this generation.
A generation that reaches $H_r$ by proposing DATA cannot append DRAIN;
it waits for commitment or recovers after Quit. Thus DRAIN is confined
to recovery from a full-quota opening. Repeated interrupted recoveries may
each add one DRAIN; there is no fixed bound across arbitrarily many such
generations. Whether some other replica has already committed the terminal
block is not an admission test.

\paragraph{Regional commitment and control history.}
Committing a physical block commits its ancestors within the regional
engine. A DRAIN commitment can therefore commit the terminal DATA ancestor.
The next region starts from the retained DATA anchor, not from a DRAIN
descendant. Old DRAIN blocks remain regional control history and evidence.
The old regional log is not rewritten when the external history selects
a shorter prefix.

\subsection{The external canonical history}
\label{sec:external}
Regional finality and external finality have different scopes.
A regional receipt proves local commitment. The external wrapper uses
the common control history to decide which regional blocks belong to
the canonical chain. Outputs of an open session are provisional at this
external level, even when they are regionally committed.
Only the retained prefix recorded by a sealed boundary is externally
irrevocable.

Write $\operatorname{End}(r)$ for the external endpoint of a closed
session. A sealed regular handoff sets $\operatorname{End}(r)=H_r$ and
retains $T_r$; activated force fixes $\operatorname{End}(r)=h_*$, even
while bootstrap is pending. For a DATA block produced in that session,
external acceptance requires
\[
 \operatorname{VerifyRegional}(B,\FC_r(B))=1,\qquad
 h(A_r)<h(B)\le\operatorname{End}(r),
\]
together with the authenticated session assignment and proof that $B$
lies on the exact DATA ancestry ending at the retained boundary anchor
(or is that anchor). This chain must extend the incoming $A_r$.
While the session is open, the same regional
checks permit only provisional output, bounded by $H_r$.

Activation installs a force cutoff immediately, even if fetching its
anchor or choosing a successor takes longer. The wrapper discards the
old session's provisional suffix above $h_*$ and waits at the retained
prefix. An external executor rebuilds its provisional state from that
prefix; the old region's BFT state is unchanged. Discarded requests
remain in the external pending pool. Late regional receipts above the
cutoff remain valid local evidence but cannot restore that external suffix.

A force must satisfy $h(A_r)\le h_*\le H_r$, so it cannot remove the
incoming sealed prefix. Once a boundary is sealed, later control records
cannot change its cut or successor. External history is therefore derived
from retained session prefixes, rather than from every regional commitment.

\subsection{Certificates, receipts, and local processing}
All messages use a chain-specific signing domain with the voting
configuration, session, region, and message type. View-dependent
messages also bind their L1 authority and $(r,g,t)$. The remaining fields
depend on the message:
\begin{itemize}
\item PROPOSE and fresh VOTE bind the opening, fresh position $j$,
block type, physical position, DATA height, and block hash.
\item NEWVIEW and ANCHOR-VOTE bind the selected opening and its block
identity. NEWVIEW also binds its justification and entry proof.
\item STATUS binds the sender's fixed entry-lock certificate and entry
proof. It does not require the next opening to have been chosen.
\item BLAME binds the admitted view only; it requires no opening, block
hash, or local progress count. All complaints for that view sign the same
statement, including when the initial leader has sent no NEWVIEW.
\end{itemize}
A proposal QC $\QC_{r,g,t}(B)$ contains matching votes of weight
$>F_{i_r}$. Correct replicas vote at most once per fresh $(r,g,t,j)$
and once for an ANCHOR per $(r,g)$. DATA and DRAIN use the same fresh
vote rule with distinct signed block types. FINAL uses the session-stable
fields defined below.

Within a session, certificate rank is $(g,t,k_r(B))$. Equal-rank choices
use a fixed total tie-break on certified content, independently of signer
subsets or encodings. An opening's identity includes its hash and positions.
The verified cross-region anchor has a distinguished base justification
$Q_r^0$ below this session's ordinary QCs. It remains an admissible
recovery fallback if no ordinary QC has formed. An unrelated region's high-ranked QC
cannot replace that anchor.

\paragraph{Regional finality receipts.}
After committing a DATA block $B$ produced in session $r$, directly or
through a descendant, a correct replica sends
\[
 \langle\tagg{FINAL},\mathit{cfg}_r,r,h(B),\Hash(B)\rangle_v
\]
to its peers and the cross-region infrastructure. Matching signatures of
weight $>F_{i_r}$ form $\FC_r(B)$. A receipt includes the block and required
ancestry, or resolves them from a validated cache. It certifies commitment,
whereas a proposal QC alone does not. It does not by itself prove
irrevocable external acceptance. If a DRAIN timer commits $T_r$, the
exported terminal receipt names $T_r$, not that DRAIN. The inherited
anchor keeps its original historical receipt.
A replica may adopt a valid receipt for its own regional session to
commit that DATA prefix locally. Receipts from other sessions serve as
anchor or history evidence; they do not replace the current engine's
locks or committed log.
The FINAL signing domain binds the chain and the displayed session-stable
fields, independently of the generation or term in which commitment
occurs. Endorsements of the same DATA block therefore combine across
recoveries; the block hash preserves its original signed context.

\paragraph{Complete validation and evidence priority.}
For $g>0$, NEWVIEW is justified by a valid QC from an older generation
of the same session, or its base justification $Q_r^0$, and its quit-entry
proof. For $g=0$, it must use the
cross-region handler's exact anchor and bootstrap proof. Its ANCHOR QC
admits the first fresh proposal. Term $t>0$ additionally requires all
preceding term-terminal QCs in that generation, each below $H_r$.
Term zero is admitted before NEWVIEW: Ready plus the boundary proof
admits $g=0$; the quit-entry proof, recovery wait, and next authority admit
$g>0$. A silent initializer therefore cannot prevent the first progress
check. Genesis is the sole synthetic initial anchor.

Validation checks signatures, authority, opening and term admission,
type and quota, both positions, parent QCs, hash links, and application
validity throughout the required ancestry. Every relevant ancestor includes
its signed proposal and block data. Correct transmissions provide complete
bundles; cached references must resolve within the complete-message bound.
Inspect signed material for independently verifiable quit evidence before
rejecting an enclosing message or discarding a duplicate. Compare it with
retained histories. At an equal timestamp, process received messages,
including evidence and complete proposals, before timer callbacks.

\paragraph{One proposal-processing rule.}
Every previously unvoted, admitted current-generation DATA or DRAIN
proposal invokes Vote, including attached history, earlier-term DATA,
and blocks already committed as ancestors. Process ancestors before
children and before terminal-QC entry. Deduplicate by
$(r,g,t,j,\Hash(B))$ after evidence inspection. Old-generation proposals
are retained but receive no new votes or timers; future-generation
proposals wait for entry. A certificate or an ANCHOR vote starts no timer.

\paragraph{Regional state and external authorization.}
The regional engine is active, frozen, or complete. Fresh proposals,
votes, and normal term transitions require an admitted, unfrozen
generation in a regionally incomplete session. A frozen generation
still accepts evidence, old votes, QCs, regional receipts, and L1 updates.
A verified terminal regional receipt completes the engine's quota and
ends local production, even if external successor selection is pending.
Local terminal commitment alone does not retire a replica: its votes
may still be needed for others' recovery until that receipt is available.

The external wrapper separately records whether a session is open,
truncated, or sealed. An activated force changes this external state,
not the regional engine's timers, locks, or voting rules. A later live
status cannot reopen the old session externally. Its continuing local
output remains subject to the installed cutoff.
The current external session is the latest unsealed session, including
a truncated session still waiting for its anchor data or successor.
Sealing its boundary advances this external cursor to the successor.
Each regional engine retains its own $(r,g,t)$ state. The external
session cursor may advance while an old engine continues. Evidence is
routed by session: it cannot freeze the successor's engine.

\begin{table}[htbp]\centering\small
\caption{Parent selection within an admitted generation.}
\label{tab:parent}
\begin{tabularx}{\linewidth}{@{}p{.32\linewidth}Y Y@{}}
\toprule
Position & Parent & Parent certificate\\
\midrule
First DATA, $t=j=0$ & $A_{r,g}$ & ANCHOR QC in $(r,g,0)$\\
First DATA, $t>0,j=0$ & Previous term's terminal DATA & Its proposal QC\\
Later DATA, $j>0$ & Position $j-1$ of this term & Its proposal QC\\
Only DRAIN, $g>0,t=j=0$ & Full-quota opening $A_{r,g}$ & ANCHOR QC in $(r,g,0)$\\
\bottomrule
\end{tabularx}
\end{table}

\section{Protocols}
\label{sec:protocols}
Protocols 1 and 2 patch v3.1's steady-state and recovery boxes. Red text
marks substantive changes from v3.1, including those retained from v4;
the session index $r$ and regional delay subscript are implicit throughout.
Protocol 3 is new relative to v3.1 and is shown in red. Explanations and service assumptions are
kept outside the execution steps.

\begin{protocol}{alg:steady}{Propose, vote, commit, and regular termination}
\begin{enumerate}[label=\textbf{\arabic*)},itemsep=5pt]
\item \textbf{Propose.} The authorized leader $\ell_{r,g,t}$ obtains the
parent and its QC from \cref{tab:parent}. \RR{If this is a DATA generation,
use the next position $j<m_{r,g,t}$. If it is a drain generation, use only
its one DRAIN position. A regionally complete session proposes nothing.}
Broadcast the complete proposal
\[
 P=\langle\tagg{PROPOSE},r,g,t,j,B,Q_{\rm parent}\rangle_{\ell_{r,g,t}}
\]
and process it locally under Vote. The parent need not be committed.

\item \textbf{Vote.} Process quit evidence first; return if it freezes
the current generation. Check completeness and validity. Request missing
data without voting; reject invalid proposals. If already voted for $P$,
stop. Otherwise forward $P$, broadcast
\RR{$\langle\tagg{VOTE},r,g,t,j,\Hash(B)\rangle_v$}, and start its
timer at the local acceptance time $\tau_v(P)$, expiring at
\RR{$\tau_v(P)+2\Delta_{i_r}$}. Process unseen earlier-term proposals
and attached history by the same rule. A duplicate never restarts a timer.

\item \textbf{Commit and export.} On timer expiry, commit $B$ and its
uncommitted ancestors. Normal term rotation preserves timers; Quit cancels
unfinished generation timers. A QC alone starts no timer.
\RR{Send FINAL endorsements for newly committed DATA produced in this
session, including the terminal DATA ancestor of a committed DRAIN.
Relay regional receipts to peers and the cross-region infrastructure;
external acceptance is determined by Protocol 3.}

\item \textbf{Terminal QC.} On a valid term-terminal proposal QC, process
its unseen history under Vote. \RR{If it certifies DATA strictly below
$H_r$}, and the generation remains unfrozen, send its complete bundle
once to $\ell_{r,g,t+1}$ and enter $(r,g,t+1)$ if not already there or
beyond it. Record first entry $e_v(r,g,t+1)$ and disable completed-term
checks. Process terminal-QC chains in order; duplicates never reset entry.
\RR{A QC for DATA at $H_r$, or for DRAIN, admits no next term or fresh
position. Continue pending timers and evidence processing.}

\item \RR{\textbf{Complete the regional engine.} On a valid $\FC_r(T_r)$,
adopt its DATA prefix, relay the terminal bundle, and mark the engine for
session $r$ complete. Stop its new proposals, votes, progress checks,
unfinished timers, and recovery. Continue serving data and running
Protocol 3. Force instructions do not invoke this step or cancel
regional timers; Protocol 3 applies their external cutoff.}
\end{enumerate}
\end{protocol}

Acceptance occurs only after complete validation and the first vote.
The term-entry clock starts at admission, before any first proposal by
that term's leader. Reaching $H_r$ in a proposal or QC is not a cross-region
entry event. With no Quit, the terminal DATA's own timer is sufficient;
no subsequent block is required for that commitment.

\paragraph{Complete quit proofs.}
For an admitted $(r,g)$, a proof is either: (a) matching
$\langle\tagg{BLAME},r,g,t\rangle$ signatures of total weight
$>F_{i_r}$ for an admitted term, including admitted term zero before
NEWVIEW; (b) two different admissible fresh
proposals signed by its leader at the same $(r,g,t,j)$; or (c) two
admissible NEWVIEWs signed by $\ell_{r,g,0}$ selecting different openings.
The proposals include their admission and quota witnesses. An excess
DATA block or unauthorized DRAIN is rejected, not treated as a second
admitted position. Different openings constitute equivocation even if
comparable; different certificates for the same opening do not.
Evidence from any admitted term can freeze its generation. Old-session
or old-generation evidence is archival and cannot freeze a current
successor. Future proofs wait for the corresponding authenticated entry.

\begin{protocol}{alg:recovery}{Blame, quit, and recovery with restricted drain}
\begin{enumerate}[label=\textbf{\arabic*)},itemsep=5pt]
\item \textbf{Blame.} At first term entry, schedule the first check at
\RR{$e_v(r,g,t)+6\Delta_{i_r}$}, even before NEWVIEW or a proposal arrives.
\RR{Once the opening is authenticated, use $q=m_{r,g,t}$ for a DATA term
and $q=1$ for a drain generation.} For $p=1,\ldots,q$, check at
\[
 e_v(r,g,t)+(2p+4)\Delta_{i_r}.
\]
If the term remains current and unfrozen and fewer than $p$ distinct
fresh positions have triggered local votes, broadcast
$\langle\tagg{BLAME},r,g,t\rangle_v$ once.
\RR{The one DRAIN vote counts; ANCHOR votes do not. Check an already-due
deadline when its opening becomes known. Do not require positions beyond $q$.}

\item \textbf{Quit.} On the first complete proof for current $(r,g)$,
freeze the generation: stop its progress checks, new proposals, votes,
and normal transitions; cancel all unfinished proposal timers.
Broadcast the proof and record quit time $a_v$ once. Retain commitments
and process old votes and QCs. Submit equivocation to L1 without awaiting
its finalization. \RR{Continue receiving regional receipts; the external watcher
processes L1 controls independently.}

\item \textbf{Status and entry.} \RR{Unless the regional engine has
completed its quota}, enter at
\[
 \RR{e_v(r,g+1,0)=a_v+\Delta_{i_r}.}
\]
\RR{The timely-authority assumption makes $\tagg{AUTH}_{r,g+1}$
available before this entry.} Read that authority. Fix the highest-ranked retained lock or valid
known QC from generations through $g$ as entry lock $Q_v$
\RR{(including the retained base justification $Q_r^0$)}.
Enter $(r,g+1,0)$, mark it unfrozen, start its progress clock, and send
$\langle\tagg{STATUS},r,g+1,0,Q_v,E\rangle_v$ with complete dependencies
to $\ell_{r,g+1,0}$. Process buffered proofs immediately. Later QCs update
retained evidence, not this fixed lock. \RR{Keep the original $A_r$,
$H_r$, and voting snapshot.}

\item \textbf{New-view.} The leader waits $2\Delta_{i_r}$ from its own
entry. If still current and unfrozen, choose the highest valid known QC
$Q^*$ from generations through $g$ \RR{or the base justification if no
ordinary QC is known}, and broadcast
\[
 \langle\tagg{NEWVIEW},r,g+1,0,A,Q^*,E\rangle_{\ell_{r,g+1,0}},
\]
where $Q^*$ certifies $A$, with complete dependencies. Choose one opening
for this generation. \RR{An opening below $H_r$ resumes DATA; an opening
at $H_r$, including a DRAIN descendant, admits exactly one fresh DRAIN.}

\item \textbf{First vote.} \RR{This handler applies to any admitted
term-zero view $(r,u,0)$, including bootstrap $u=0$.} Process quit evidence,
then validate NEWVIEW and its dependencies. If current and unfrozen,
no ANCHOR vote was cast, and $\rank(Q^*)\ge\rank(Q_v)$, forward NEWVIEW
and broadcast
$\langle\tagg{ANCHOR\mbox{-}VOTE},r,u,0,\Hash(A)\rangle_v$.
Start no timer and consume no fresh position. On an ANCHOR QC
\RR{of weight $>F_{i_r}$}, the leader resumes Protocol 1 at $j=0$.
\end{enumerate}
\end{protocol}

The entry-lock test governs ANCHOR voting only. A subsequent fresh
proposal backed by a valid ANCHOR QC is processed under Protocol 1 even
at a replica that declined the ANCHOR vote. A locally committed terminal
block does not justify stopping this recovery before a terminal regional
receipt is available. An external force cutoff is not a local quit proof
and does not alter this recovery rule.

\subsection{Boundary witnesses and the bootstrap event}
\label{sec:boundary}
For an effective map $\chi$, let $\tagg{NextLive}(i,\chi)$ be the first
other live region after $i$ in $\mathcal L_R$, searching at most one cycle.
If none exists, its value is $\bot$. A boundary witness has the form
\[
 \mathcal W=(\mathit{mode},r,r+1,i',A,\FC(A),\lambda_s,\Xi),
 \qquad \mathit{mode}\in\{\tagg{REGULAR},\tagg{FORCE}\},
\]
where $\Xi$ authenticates the control order, effective map at selection
height $\lambda_s$, and any activated L1 force entry. It also
binds the successor authority and configuration. The anchor's hash is
fixed by the block and finality receipt inside $\mathcal W$, not by the
multisignature height tag. The accepted witness seals that successor,
retained DATA prefix, and selection height. Certificate encodings do not
change boundary identity. A force witness requires no evidence that the
old region has stopped or lacks higher regional commitments.
Every external verifier can validate this witness and install its sealed
prefix; membership in the successor region is needed only to bootstrap
that region's engine.

For a regular boundary, $A=T_r$, its ancestry contains exactly $TB_r$
DATA blocks since $A_r$, and selection starts at the terminal block's
signed L1 height. For a forced boundary, $A$ is the regionally committed DATA block
at $h_*$ named by the activated entry, with
$h(A_r)\le h_*\le H_r$. If $h_*=h(A_r)$, reuse the incoming anchor and
its historical receipt. No new receipt from the failed region is required
for a prefix it inherited unchanged.
Forced selection uses $\lambda_s\ge\max\{a_j,\lambda(A)\}$ and replays
all due controls through $\lambda_s$ before choosing the successor.
For genesis, the initial L1 reference supplies $\lambda(A_0)$.

At an unsealed boundary, process due controls in their authenticated
order and apply an applicable force before regular quota completion at
the same control height. A sealed boundary is immutable. A force targeting
an already departed session cannot replace its successor; its status map
still applies in the ordered control history. A later live update does
not erase the installed cutoff or reopen that session's external
authorization.
If $i'=\bot$, retain the anchor and advance the authenticated control
cursor; the first admissible later prefix with an eligible successor seals
the waiting boundary. A newer local L1 head cannot reselect it afterward.

\begin{definition}[Local cross-region bootstrap event]
\label{def:bootstrap}
For a replica $v\in V_{i'}$, $\tagg{Ready}_v(\mathcal W)$ occurs on its
first complete validation of an admissible sealed witness naming $i'$,
the anchor and missing application history, and the matching finalized
$\tagg{AUTH}_{r+1,0}$. REGULAR requires terminal finality and the regional
quota; FORCE additionally requires finalized L1 through its activation
height and the corresponding external cutoff in the ordered control history. The event is suppressed if this replica
already installed that successor or has authenticated a later session.
An incomplete boundary stays pending until its data and authority arrive.
\end{definition}
This is the precise event at which a node recognizes that its region may
start the next session. Merely receiving a leading block at height $H_r$,
its QC, an announced force entry, or a newer local status map is insufficient.
Ready starts the local entry clock; a remote proposal time, private QC
assembly, or L1 installation time does not.

\paragraph{Terminal recovery example.}
Let $h(A_r)=100$, $TB_r=3$, and $K=2$. The first term proposes DATA
101--102 and its terminal QC admits the next term, which has only DATA
103. Its QC starts no next term. If its own timer survives, the receipt
for DATA 103 enables the regular handler. If Quit cancels the timers and
the next generation opens on certified DATA 103, its sole fresh proposal
is one DRAIN. That DRAIN's timer commits DATA 103 as an ancestor and
produces the same terminal DATA receipt. Another Quit may require a new
generation and one further DRAIN. In every case, the next region starts
with DATA 104, whose parent is DATA 103.

\paragraph{Forced bootstrap example.}
Suppose a session opened at DATA 100 with $TB_r=10$. A multisignature
entry specifies cutoff 106 and force activation at L1 height 1200.
The old region may locally commit DATA 107 and 108 during the notice
period or afterward. Once L1 through 1200 is finalized, external nodes
retain this session only through DATA 106. They discard any provisional
old-session suffix above 106 and recheck its requests for inclusion.
After verifying DATA 106 and its regional receipt, the selected successor
opens $(r+1,0,0)$ and proposes its own DATA 107 with parent DATA 106.
The old regional log is unchanged. Its receipts for 107 and 108 do not
restore those blocks to the external canonical chain.

\begin{protocol}{alg:bootstrap}{Watch the stream and bootstrap the next region}
\color{rrRed}
\begin{enumerate}[label=\textbf{\arabic*)},itemsep=5pt]
\item \textbf{Watch and fetch.} Every prospective successor continuously
tracks the externally active session. Receive leading blocks and QCs,
fetch missing ranges and regional receipts, and independently poll
finalized L1 for authorities and status updates. Validate and cache
dependencies. Apply the ordered controls before accepting external
output: exclude any session suffix above its installed cutoff, even if
its block carries an older L1 reference. Open-session output is
provisional; it starts no regional progress clock.
On a valid sealed boundary witness, install its retained DATA prefix,
endpoint, and successor in the external state, and update the pending
pool. This external update does not require membership in the successor.

\item \textbf{Handle regular completion.} On a terminal DATA bundle,
verify its regional receipt and exactly $TB_r$ DATA positions. Starting
from its signed L1 reference, replay the ordered boundary controls.
If no applicable force has taken precedence, derive
$i'=\tagg{NextLive}(i_r,\chi_{\rm eff}(\lambda_s))$ and verify the
regular witness retaining the prefix through $H_r$. If none is eligible,
keep the candidate anchor and continue watching controls. A terminal
regional receipt cannot override an installed force cutoff.

\item \textbf{Handle activated force.} On finalized L1 through $a_j$,
validate $\tagg{FORCE}(r,h_*)$ against the targeted externally active
session and install $\operatorname{End}(r)=h_*$ once. Discard its
provisional external suffix above $h_*$; restore the retained application
state when its data are available, and keep orphaned admissible requests
pending. Do not alter old regional timers, locks, or commitments.
Fetch $B_{h_*}$ with its historical regional receipt and ancestry.
Replay controls through $\lambda_s\ge\max\{a_j,\lambda(B_{h_*})\}$ and
derive the successor. If data or an eligible region is missing, keep the
cutoff and wait. A stale target cannot reopen a sealed boundary.
No old-region shutdown proof or fresh response is required.

\item \textbf{Enter on Ready.} When $\tagg{Ready}_v(\mathcal W)$ fires
for this replica's region, relay the complete bundle within that region.
Install session $r+1$, its exact anchor $A_{r+1}=A$, configuration, quota,
and opening control height $\lambda_s$. Record the boundary's retained
DATA prefix through the external update in step 1. Set its retained and entry lock
$Q_v$ to the fixed base justification $Q_{r+1}^0$. Enter $(r+1,0,0)$, record
$e_v(r+1,0,0)$ once, mark it unfrozen, and start the first progress check
before any leader message. Send STATUS containing the fixed base anchor
justification and $\mathcal W$ to $\ell_{r+1,0,0}$. Process buffered
current-generation evidence. A duplicate never resets entry, quota, or locks.

\item \textbf{Announce and certify the opening.} After $2\Delta_{i'}$
from its entry, the initial leader, if still current and unfrozen,
broadcasts NEWVIEW selecting exactly $A_{r+1}$, with
$Q^*=Q_{r+1}^0$ and $\mathcal W$.
Members validate the fixed anchor, authority, and entry proof, then forward
NEWVIEW and cast their one ANCHOR vote under Protocol 2's first-vote rule.
ANCHOR voting starts no timer. No highest-QC choice from the old region
may substitute another block for this opening.

\item \textbf{Start the region.} On the new region's ANCHOR QC of weight
$>F_{i'}$, its initial leader invokes Protocol 1 for the first DATA block
at global height $h(A_{r+1})+1$ and local physical position $1$.
Use the new session's fresh $TB_{r+1}$ quota. Subsequent local failures
invoke Protocol 2 within this session; they do not themselves increment $r$.
Continue watching L1 and streaming the new region's blocks and receipts.
\end{enumerate}
\end{protocol}

Re-run both handlers on relevant stream, authority, or control updates,
including while regional consensus is frozen. They are idempotent.
Before installing or producing under a successor, replay further
authenticated controls already due for it. The cutoff remains installed
while bootstrap waits. Genesis uses a synthetic initial boundary and
the same bootstrap exchange. After Ready, even a silent initial leader
faces the first check at $e_v(r+1,0,0)+6\Delta_{i'}$.

\begin{thebibliography}{9}
\bibitem{sync}
I. Abraham, D. Malkhi, K. Nayak, L. Ren, and M. Yin.
\emph{Sync HotStuff: Simple and Practical Synchronous State Machine Replication}.
IEEE Symposium on Security and Privacy, 2020, pp. 106--118.
\href{https://conferences.computer.org/sp/pdfs/sp/2020/349700a106.pdf}{Published paper}.
\end{thebibliography}
\end{document}
